\documentclass[12pt]{article}

\usepackage[portuguese]{babel}
\usepackage{float}
\usepackage{listings}
\usepackage{amsfonts, amssymb, amsmath}
\usepackage{graphicx}
\usepackage{enumitem}
\usepackage[colorlinks=true, allcolors=blue]{hyperref}
\usepackage[letterpaper,
            top=1.5cm,
            bottom=1.5cm,
            left=2.75cm,
            right=2.75cm,
            marginparwidth=1.75cm]{geometry}

\title{Programação Funcional}
\author{Prof. Malbarbo}
\date{\today}

\begin{document}
\maketitle
\tableofcontents
\vspace{48pt}

Dentro da programação não temos apenas uma forma de implementar um programa, 
temos alguns diferentes paradigmas que podemos usar:
\begin{itemize}[label=$\rightarrow$]
    \item Programação Imperativa -- Os programas são descritos com sentenças 
    que modificam o estado do programa
    \item Programação Orientada a Objetos -- Os programas são descritos 
    organizados em torno de objetos, classes e interações
    \item Programação Funcional -- Os programas são descritos com aplicação 
    e composição de funções
    \begin{itemize}
        \item Evita mudança de estado (mudança do valor das variáveis)
        \item Evita efeitos colaterais, que são quaisquer efeitos que sejam 
        observáveis além do valor de saída da função, como a mudança dos 
        parâmetros e variáveis globais, exceções, entrada e saída, etc.
    \end{itemize} 
\end{itemize}

\section{Por que usar o paradigma funcional?}
Um paradigma de programação é uma ferramenta, e conhecer várias ferramentas 
permite utilizar a mais adequada para cada problema, e muitas vezes o Python e 
a programação imperativa têm problemas.
\begin{itemize}
    \item Mudança de estados:
    \begin{itemize}
        \item Devido a forma com que uma lista é alocada em Python, 
        inconsistencias acontecem, como quando alocamos:
        \begin{verbatim}
    >>> lst = [0] * 3
    >>> lst
    [0, 0, 0]
    >>> lst[1] = 10
    >>> lst
    [0, 10, 0]
        \end{verbatim}
        \item Porem, ao alocar uma lista dentro de lista:
        \begin{verbatim}
    >>> lst = [[]] * 3
    >>> lst
    [[], [], []]
    >>> lst[1].append(2)
    >>> lst
    [[2], [2], [2]]
        \end{verbatim}
        \item Além disso, quando modificando uma lista usando 
        \verb|adiciona_todos|
        \begin{verbatim}
def adiciona_todos(dest: list[int], fonte: list[int]):
    ''' 
    Adiciona todos os elementos de *fonte* no final de 
    *dest* 
    '''
    for x in fonte:
        dest.append(x)
        \end{verbatim}
        \begin{itemize}
            \item Ao modificar lst usando uma lista diferente não temos 
            problemas
            \begin{verbatim}
    >>> lst = [4, 3, 1]
    >>> adiciona_todos(lst, [6, 2])
    >>> lst
    [4, 3, 1, 6, 2]
            \end{verbatim}
            \item Porém, ao adicionar lst nela mesma:
            \begin{verbatim}
    >>> adiciona_todos(lst, lst)
    >>> lst
            \end{verbatim}
            A execução não acaba, visto que a lista se automodifica 
            infinitamente.
        \end{itemize}
    \end{itemize}
    \item Efeitos colaterais
    \begin{itemize}
        \item Dentro de uma função, como vimos, é possível mudar o valor de 
        um dos parâmetros de entrada, dependendo da alteração, o mesmo 
        código, dependendo da ordem que são realizadas as funções, temos 
        resultados diferentes.
        \begin{verbatim}
def soma_indices(lst: list[int], a: int, b: int) -> int:
    return indice(lst, b) + indice(lst, a)

def soma_indices(lst: list[int], a: int, b: int) -> int:
    return indice(lst, a) + indice(lst, b)
        \end{verbatim}
        Não é possível que as duas definições são equivalentes sem olhar o 
        código da função \verb|indice|. Se ela tiver efeitos colaterais, 
        então as definições não podem ser equivalentes.
    \end{itemize}
    \item Com os efeitos colaterias dificulta pensar localmente sobre o 
    funcionamento do código, sem eles, podemos fazer isso.
\end{itemize}

\section{Fundamentos da Programação Funcional}
O paradigma de programação funcional é baseado na definição e aplicação de funções
\begin{itemize}
    \item Função: Conjunto de expressões que mapeia valores de entrada para 
    valores de saída
    \item Expressão: Entidade sintática que quando avaliada (reduzida) produz um 
    valor
    \item Sentença: Operação que não retorna valor, apenas um efeito colateral 
    (no Python, if, for, while, etc) -- Não temos sentenças no paradigma funcional
\end{itemize}
Uma expressão consiste em:
\begin{itemize}
    \item Um literal -- Valor que é representado diretamente no código, em geral, 
    os literais utilizados para criar valores de tipos primitivos.
    \item Uma função primitiva -- Função suportada diretamente pela linguagem 
    de programação
\end{itemize}
Tipos primitivos:
\begin{itemize}
    \item No Gleam, os tipos primitivos sempre começam com letra maiúscula
    \item Numero inteiro (Int)
    \begin{itemize}
        \item \verb|1345|
        \item \verb|9_876|
    \end{itemize}
    \item Numero com ponto flutuante (Float)
    \begin{itemize}
        \item \verb|2.65|
        \item \verb|2.0e12| 
        \item \verb|7.4e-10|
    \end{itemize}
    \item Booleano (Bool)
    \begin{itemize}
        \item True
        \item False
    \end{itemize}
    \item Strings (String)
    \begin{itemize}
        \item \verb|"din uem"|
        \item \verb|"apenas \"um\" teste"|
    \end{itemize}
\end{itemize}
Funções primitivas:
\begin{itemize}
    \item Operações com inteiros:
    \begin{itemize}
        \item \verb|+ (int.add)|
        \item \verb|- (int.subtract)|
        \item \verb|* / % > >= <= == !=| 
    \end{itemize}
\end{itemize}

\end{document}


